\begin{proposition}[Maximal monodromy with trivial central dualizing bundle]\hypertarget{tag-030X}{}\label{push:src:B.4.28}
\StatementRecord{push:src:B.4.28}{proved}{030X}{\citet[Definition 1.1 and Lemmas 1.2, 2.2, 2.7]{Friedman2016AB}; \citet[Corollary 2.25, Lemma 2.26 and Proposition
2.27]{Dloussky2005AB}; \citet[Sections 3.2--3.3]{FujikiPontecorvo2010AB}; \citet[Section 15, Theorem 15.1 and the following
paragraph, pp.~130--131]{Oda1978D}; \citet[Chapter~III, Theorem~4.12(ii), pp.~134--135]{BS76}; \citep[Chapter~III, Theorem~4.12(iii), pp.~134--135]{BS76}}

Let $p:\mathcal X\to\Delta$ be a proper holomorphic map
from a smooth complex threefold, smooth over $\Delta^*$,
whose nonzero fibres are complex two-tori. Suppose that
$D=p^{-1}(0)$ is a reduced normal-crossing divisor and that $\omega_D\simeq\OO_D$.
If its monodromy on $H_1(\mathcal X_t,\Q)$ satisfies
\begin{equation}
 \operatorname{rank}(T-1)=2,
\label{pub:I:eq:B-4-28-1}
\end{equation}
then every normalization component of $D$ is a smooth
rational toric surface, its conductor is its toric
boundary, and the dual complex of $D$ is a topological
two-torus. In particular all normalized components are
K\"ahler.

The trivialization of $\omega_D$ in the hypotheses is the central
section used in the residue and nearby-period arguments below.
\end{proposition}


\begin{proof}
\begin{enumerate}[label=\arabic*.,leftmargin=*]
\item Write $V_i$ for the compact smooth normalization
components and $C_i$ for their reduced conductor divisors.
Adjunction gives
\begin{equation}
 K_{V_i}+C_i\sim0.
\label{pub:I:eq:B-4-28-2}
\end{equation}
If $D$ is smooth, the map is a proper submersion near
its compact central fibre and its monodromy is trivial.
Otherwise every $C_i$ is nonempty, by connectedness.
Adjunction on a conductor branch says that it is either
a smooth elliptic curve disjoint from the other branches,
or a rational curve with two boundary points, counted
on its normalization. The rational branches form cycles,
including the nodal irreducible and two-component cases.

\item \emph{Nonnegative Euler contributions.}
A nonempty effective anticanonical divisor makes every
pluricanonical system empty: multiplying a nonzero
pluricanonical section by the appropriate power of the
anticanonical section would give a nonzero holomorphic
function with a nonempty zero divisor. Thus the surface
has Kodaira dimension $-\infty$. Its minimal model is
rational, ruled, or of class VII. We spell out the
consequences for the reduced anticanonical pairs here.

\item On a rational surface, the conductor is connected, by
$0\to K_V\to\OO_V\to\OO_C\to0$ and $H^1(K_V)=0$.
For a rational cycle, the charge identity and positivity
are
\begin{equation}
 e(V\setminus C)=12-C^2-\#\{\text{cycle components}\}
 \ge0,\qquad
 e(V\setminus C)=0\ \Longleftrightarrow\ (V,C)
 \text{ is toric}.
\label{pub:I:eq:B-4-28-3}
\end{equation}
Here $e$ denotes compactly supported Euler characteristic.
This includes the nodal irreducible convention:
\citet[arXiv version~2, Section~1, p.~3; Definition~1.1 and Lemma~1.2, pp.~4--5; Lemma~2.2, p.~11; Lemma~2.7, p.~14]{Friedman2016AB}.
For a smooth elliptic anticanonical curve,
$e(V\setminus C)=2+b_2(V)>0$ instead.

\item On a nonrational ruled minimal model, a conductor
component dominating the base has degree at most two
and genus at most one. The base is therefore elliptic.
The minimal conductor is two disjoint elliptic sections
or one unramified elliptic bisection. Its Euler
contribution is zero. The two section squares sum to
zero, and the bisection has square zero, also directly
from $K_V+C\sim0$ and $K_V^2=0$.

\item A minimal class-VII surface with $b_2=0$ and a curve is
a Hopf surface. Its elliptic curves have square zero,
and a reduced anticanonical divisor has at most two
elliptic components. For this last assertion the exact
sequence $0\to K_V\to\OO_V\to\OO_C\to0$ gives
$h^0(\OO_C)\le1+h^1(K_V)=2$.
The Euler contribution is zero. This also covers
secondary Hopf surfaces.

\item For a minimal class-VII surface with $b_2>0$, the
reduced anticanonical divisor is a numerically
anticanonical divisor whose coefficients are all one.
\citet[arXiv version~3, Corollary~2.25, pp.~16--17; Lemma~2.26, p.~17; Proposition~2.27, pp.~17--19]{Dloussky2005AB} imply that it is parabolic Inoue
or Inoue--Hirzebruch. Half-Inoue has only a nontrivially
twisted anticanonical divisor and cannot satisfy
\eqref{pub:I:eq:B-4-28-2}. The remaining possibilities are hyperbolic
Inoue with two disjoint rational cycles, or parabolic
Inoue with one rational cycle and a disjoint elliptic
curve. If $n=b_2>0$, there are $n$ rational cycle
components and $e(V)=n$, so the contribution is zero.
In the parabolic case the elliptic curve has square
$-n$. The curve configurations, anticanonical identities
and underlying $(S^1\times S^3)\#n\overline{\PP}^2$
topology are recorded in
\citet[Sections~3.2--3.3, pp.~20--22]{FujikiPontecorvo2010AB}.

\item These assertions extend to nonminimal pairs. An
exceptional curve outside $C$ meets it once, transversely
at a smooth point, since $-K_V\cdot E=1$; its contraction
reverses an interior blow-up. An exceptional curve in
$C$ meets the rest at two points and its contraction
reverses a corner blow-up. The reduced anticanonical
condition is preserved. Interior blow-up increases
$e(V\setminus C)$ by one, whereas corner blow-up leaves
it unchanged. Thus every contribution is nonnegative.
Equality permits only toric rational pairs, minimal
elliptic ruled or Hopf pairs, and corner blow-ups of
hyperbolic or parabolic Inoue pairs. For these Inoue
pairs the cycle count and $b_2$ increase together, and
the fundamental group stays $\Z$. Corner blow-ups do
not touch the parabolic elliptic curve, whose square
remains strictly negative.

\item The local nearby fibre at a multiplicity-$k$ stratum
of a reduced normal-crossing divisor retracts to
$(S^1)^{k-1}$. Euler integration consequently gives
\begin{equation}
 0=e(\mathcal X_t)=\sum_i e(V_i\setminus C_i).
\label{pub:I:eq:B-4-28-4}
\end{equation}
Every summand is nonnegative and therefore zero.
Rational components with a smooth elliptic conductor,
and all interior blow-ups, are excluded.

\item \emph{Elliptic normal-degree balance.}
An elliptic double curve contains no triple point.
The local model $p=xy$ gives opposite degrees to its
two normal lines. Sum these degrees over all elliptic
conductor branches. The total is zero. An elliptic
ruled or Hopf component contributes zero; a parabolic
Inoue component contributes the strictly negative
degree of its sole elliptic branch. No rational
component with elliptic conductor remains. Hence there
are no parabolic Inoue components.

\item By connectedness, all conductors are now elliptic or
all are rational: no remaining component has both
types. In the elliptic case there are no triple points
and every vertex of the dual graph has valence at most
two. The graph is a segment or circle, allowing a loop
or multiple edges. In this case
\begin{equation}
 \operatorname{rank}(T-1)\le b_1(\text{dual graph})
 \le1.
\label{pub:I:eq:B-4-28-5}
\end{equation}
To prove this, remove
small tubes about the disjoint double curves. The
remaining compact cores extend as smooth submersions
over the whole disc and can be trivialized there.
Monodromy is then a product of the annular twists
$xy=t$. A passage of a loop across one neck adds its
oriented normal-circle class, with coefficient the
signed crossing number of that edge. Core loops have
zero change. Crossing numbers of a closed loop form a
cycle of the dual graph. Thus $T-1$ factors through its
rational cycle space. The normal-circle rotation exists
on the circle bundle of the normal line. This proves
\eqref{pub:I:eq:B-4-28-5} and excludes the elliptic case.

\item \emph{Hyperbolic components contradict a nonzero nearby period.}
The remaining possibilities are toric rational pairs
and corner blow-ups of hyperbolic Inoue pairs. Suppose
one component is of the latter kind. A boundary of a
small neighborhood of either rational conductor cycle,
lying in $V_i\setminus C_i$, is a torus mapping torus
$M_A$, whose monodromy on its fibre lattice is
\begin{equation}
 A\in\operatorname{SL}_2(\Z),\qquad
 |\operatorname{tr}A|>2.
\label{pub:I:eq:B-4-28-6}
\end{equation}
Here is the exact marking needed. In
\citet[Section 15, Theorem 15.1 and the following
paragraph, retypeset pp.~130--131]{Oda1978D}, the universal
toric construction has lattice $N$ and deck action
given by multiplication by a nontrivial totally
positive unit of a real quadratic field. This is
the hyperbolic matrix $A$. The complement of the
two cycles is $\R$ times the compact angular-torus
bundle over a circle. A path in the logarithmic
absolute-value space from a point to its image by
$A$ realizes this mapping torus. The point can be
chosen in a node cone arbitrarily near a conductor
node. Its two adjacent primitive rays form a lattice
basis, so the angular fibre is exactly the Clifford
torus in the corresponding local coordinates
$(x,y)$. Corner blow-ups are toric star subdivisions;
they alter neither the dense torus nor $A$, only
this unimodular coordinate basis. Thus the same
marked mapping torus exists for every corner
blow-up allowed above.

\item Choose a compact such $M_A$ in the smooth open stratum.
A neighborhood of this compact set extends as a
smooth submersion over the full parameter disc.
Transport gives an embedding $M_A\to\mathcal X_t$ for
all small $t$, with the same marked fibre $L_t\simeq
(S^1)^2$. In the abelian group
$H_1(\mathcal X_t,\Q)$, the mapping-torus relations give
\[
 (A-1)H_1(L_t,\Q)\longmapsto0.
\]
Since $A-1$ is invertible over $\Q$, both angular
generators map to zero. The cohomology ring of a
complex two-torus is the exterior algebra on $H^1$.
It follows that the map $H^2(\mathcal X_t,\Q)\to
H^2(L_t,\Q)$ is zero and hence
\begin{equation}
 [L_t]=0\quad\hbox{in }H_2(\mathcal X_t,\Q).
\label{pub:I:eq:B-4-28-7}
\end{equation}

\item On the other hand, $h^0(D,\omega_D)=1$, since $D$ is
compact and connected, and the corresponding dimension
on every nearby two-torus is also one. Grauert's
cohomology and base-change theorem applied to the
invertible relative dualizing sheaf supplies a relative
section extending the chosen nonzero central generator;
see \citet[Chapter~III, Theorem~4.12(ii), pp.~134--135]{BS76}.
Evaluation is
nonzero at every point of $D$. Compactness and shrinking
make this section nonvanishing near the entire fibre.
Write it as $\Omega$. At the conductor node just used,
the total family has coordinates $xyz=t$ and
\begin{equation}
 \Omega_t=U(x,y,z)\frac{dx}{x}\wedge\frac{dy}{y},
 \qquad U(0,0,0)\ne0.
\label{pub:I:eq:B-4-28-8}
\end{equation}
The angular torus can be transported specifically as
$|x|=\epsilon_1$, $|y|=\epsilon_2$, $z=t/(xy)$.
This agrees, up to isotopy, with its compact-core
transport. If $U=\sum a_{ijk}x^iy^jz^k$ in a small
polydisc, direct integration of the convergent series
gives
\begin{equation}
 \int_{L_t}\Omega_t
   =(2\pi i)^2\sum_{k\ge0}a_{kkk}t^k
   \longrightarrow(2\pi i)^2U(0,0,0)\ne0.
\label{pub:I:eq:B-4-28-9}
\end{equation}
For small $t$ this contradicts \eqref{pub:I:eq:B-4-28-7}, because
the holomorphic two-form $\Omega_t$ is closed. This
excludes every hyperbolic Inoue component.

\item \emph{The dual complex.}
All normalization pairs are now rational toric pairs.
Every vertex link is one circle and every edge has
two incident faces, because each rational double
curve has two triple points. Thus the finite dual
complex $K$ is a connected closed topological surface.
It is oriented: the iterated residues of a
trivialization of $\omega_D$ are nonzero at every
triple point, and the residue theorem on each rational
double curve makes its two coefficients opposite,
with the gluing sign for the two branches. These
coefficients form a nowhere-zero two-cycle on $K$.
Constancy of coherent Euler characteristic in the
proper flat analytic family
\citep[Chapter~III, Theorem~4.12(iii), pp.~134--135]{BS76}, and the alternating
normalization resolution, give
\begin{equation}
 0=\chi(\OO_D)=\chi(K).
\label{pub:I:eq:B-4-28-10}
\end{equation}
The surface therefore has genus one.

\item For self-crossings, use the finite local-branch
incidence cell structure; the link of a normalized
toric component is its circular boundary. Over the
normalization $\PP^1$ of a double curve, the two-branch
cover extends unramified over its triple-point ends
by local normal-crossing coordinates and hence splits.
Subdivision
displays the same closed surface and the residue signs.
The alternating normalization resolution uses these
branch signs and gives the same Euler count. The
hyperbolic-link argument took place in the open
normalization stratum, which embeds in $D$, and is
unchanged by self-identification of boundary branches.
\end{enumerate}
\end{proof}